% figure4.tex  –  TikZ diagram: CNN architecture for histopathology
% Included via % figure4.tex  –  TikZ diagram: CNN architecture for histopathology
% Included via % figure4.tex  –  TikZ diagram: CNN architecture for histopathology
% Included via % figure4.tex  –  TikZ diagram: CNN architecture for histopathology
% Included via \input{../figures/figure4.tex} from chapter4.tex
% Requires: tikz, tikz libraries (shapes.geometric, arrows.meta, positioning, fit, backgrounds)

\begin{tikzpicture}[
    font=\small,
    >=Stealth,
    node distance=0.55cm and 0.9cm,
    % ---- block styles ----
    imgblock/.style={
        draw, fill=pink!40, rounded corners=2pt,
        minimum width=0.7cm, minimum height=1.6cm,
        align=center
    },
    convblock/.style={
        draw, fill=blue!25, rounded corners=2pt,
        minimum width=0.55cm, minimum height=1.4cm,
        align=center
    },
    poolblock/.style={
        draw, fill=green!30, rounded corners=2pt,
        minimum width=0.55cm, minimum height=1.0cm,
        align=center
    },
    fcblock/.style={
        draw, fill=orange!30, rounded corners=2pt,
        minimum width=0.55cm, minimum height=0.9cm,
        align=center
    },
    outblock/.style={
        draw, fill=red!25, rounded corners=4pt,
        minimum width=1.1cm, minimum height=0.7cm,
        align=center
    },
    lbl/.style={font=\scriptsize, align=center},
    arrow/.style={->, thick, gray!70}
]

%% ---- Input image ----
\node[imgblock] (img) {};
\node[lbl, below=0.15cm of img] {Input\\image\\patch};

%% ---- Conv block 1 (3 feature maps shown as stacked rectangles) ----
\foreach \i in {0,1,2}{
    \node[convblock, right=\i*0.12cm+0.9cm of img, yshift=\i*0.12cm] (conv1\i) {};
}
\node[lbl, below=0.15cm of conv10] {Conv 1\\(ReLU)};

%% ---- Pool 1 ----
\foreach \i in {0,1,2}{
    \node[poolblock, right=\i*0.12cm+0.55cm of conv10, yshift=\i*0.12cm] (pool1\i) {};
}
\node[lbl, below=0.15cm of pool10] {Pool 1\\(Max)};

%% ---- Conv block 2 ----
\foreach \i in {0,1,2,3}{
    \node[convblock, right=\i*0.12cm+0.55cm of pool10, yshift=\i*0.12cm] (conv2\i) {};
}
\node[lbl, below=0.15cm of conv20] {Conv 2\\(BN+ReLU)};

%% ---- Pool 2 ----
\foreach \i in {0,1,2,3}{
    \node[poolblock, right=\i*0.12cm+0.55cm of conv20, yshift=\i*0.12cm] (pool2\i) {};
}
\node[lbl, below=0.15cm of pool20] {Pool 2\\(Max)};

%% ---- Conv block 3 ----
\foreach \i in {0,...,4}{
    \node[convblock, right=\i*0.12cm+0.55cm of pool20, yshift=\i*0.12cm] (conv3\i) {};
}
\node[lbl, below=0.15cm of conv30] {Conv 3\\(BN+ReLU)};

%% ---- Flatten arrow label ----
\node[right=0.65cm of conv34, yshift=-0.3cm] (flat) {};

%% ---- FC layers ----
\node[fcblock, right=1.5cm of conv30] (fc1) {};
\node[lbl, below=0.15cm of fc1] {FC 1\\(Dropout)};

\node[fcblock, right=0.7cm of fc1] (fc2) {};
\node[lbl, below=0.15cm of fc2] {FC 2\\(Dropout)};

%% ---- Output ----
\node[outblock, right=0.7cm of fc2] (out) {Softmax\\output};
\node[lbl, below=0.15cm of out] {Class\\probs.};

%% ---- Arrows ----
\draw[arrow] (img.east)    -- (conv10.west);
\draw[arrow] (conv12.east) -- (pool10.west);
\draw[arrow] (pool12.east) -- (conv20.west);
\draw[arrow] (conv23.east) -- (pool20.west);
\draw[arrow] (pool23.east) -- (conv30.west);
\draw[arrow] (conv34.east) -- node[above, font=\scriptsize]{Flatten} (fc1.west);
\draw[arrow] (fc1.east)    -- (fc2.west);
\draw[arrow] (fc2.east)    -- (out.west);

%% ---- Brace annotations above ----
\draw[decorate, decoration={brace, amplitude=5pt, raise=4pt}, thick, blue!60]
    (conv10.north west) -- (pool12.north east)
    node[midway, above=8pt, font=\scriptsize, blue!80] {Feature extraction stage 1};

\draw[decorate, decoration={brace, amplitude=5pt, raise=4pt}, thick, blue!60]
    (conv20.north west) -- (pool23.north east)
    node[midway, above=12pt, font=\scriptsize, blue!80] {Feature extraction stage 2};

\draw[decorate, decoration={brace, amplitude=5pt, raise=3pt}, thick, blue!60]
    (conv30.north west) -- (conv34.north east)
    node[midway, above=14pt, font=\scriptsize, blue!80] {Deep feature maps};

\draw[decorate, decoration={brace, amplitude=5pt, raise=3pt}, thick, orange!70]
    (fc1.north west) -- (fc2.north east)
    node[midway, above=10pt, font=\scriptsize, orange!90] {Classification head};

\end{tikzpicture}
 from chapter4.tex
% Requires: tikz, tikz libraries (shapes.geometric, arrows.meta, positioning, fit, backgrounds)

\begin{tikzpicture}[
    font=\small,
    >=Stealth,
    node distance=0.55cm and 0.9cm,
    % ---- block styles ----
    imgblock/.style={
        draw, fill=pink!40, rounded corners=2pt,
        minimum width=0.7cm, minimum height=1.6cm,
        align=center
    },
    convblock/.style={
        draw, fill=blue!25, rounded corners=2pt,
        minimum width=0.55cm, minimum height=1.4cm,
        align=center
    },
    poolblock/.style={
        draw, fill=green!30, rounded corners=2pt,
        minimum width=0.55cm, minimum height=1.0cm,
        align=center
    },
    fcblock/.style={
        draw, fill=orange!30, rounded corners=2pt,
        minimum width=0.55cm, minimum height=0.9cm,
        align=center
    },
    outblock/.style={
        draw, fill=red!25, rounded corners=4pt,
        minimum width=1.1cm, minimum height=0.7cm,
        align=center
    },
    lbl/.style={font=\scriptsize, align=center},
    arrow/.style={->, thick, gray!70}
]

%% ---- Input image ----
\node[imgblock] (img) {};
\node[lbl, below=0.15cm of img] {Input\\image\\patch};

%% ---- Conv block 1 (3 feature maps shown as stacked rectangles) ----
\foreach \i in {0,1,2}{
    \node[convblock, right=\i*0.12cm+0.9cm of img, yshift=\i*0.12cm] (conv1\i) {};
}
\node[lbl, below=0.15cm of conv10] {Conv 1\\(ReLU)};

%% ---- Pool 1 ----
\foreach \i in {0,1,2}{
    \node[poolblock, right=\i*0.12cm+0.55cm of conv10, yshift=\i*0.12cm] (pool1\i) {};
}
\node[lbl, below=0.15cm of pool10] {Pool 1\\(Max)};

%% ---- Conv block 2 ----
\foreach \i in {0,1,2,3}{
    \node[convblock, right=\i*0.12cm+0.55cm of pool10, yshift=\i*0.12cm] (conv2\i) {};
}
\node[lbl, below=0.15cm of conv20] {Conv 2\\(BN+ReLU)};

%% ---- Pool 2 ----
\foreach \i in {0,1,2,3}{
    \node[poolblock, right=\i*0.12cm+0.55cm of conv20, yshift=\i*0.12cm] (pool2\i) {};
}
\node[lbl, below=0.15cm of pool20] {Pool 2\\(Max)};

%% ---- Conv block 3 ----
\foreach \i in {0,...,4}{
    \node[convblock, right=\i*0.12cm+0.55cm of pool20, yshift=\i*0.12cm] (conv3\i) {};
}
\node[lbl, below=0.15cm of conv30] {Conv 3\\(BN+ReLU)};

%% ---- Flatten arrow label ----
\node[right=0.65cm of conv34, yshift=-0.3cm] (flat) {};

%% ---- FC layers ----
\node[fcblock, right=1.5cm of conv30] (fc1) {};
\node[lbl, below=0.15cm of fc1] {FC 1\\(Dropout)};

\node[fcblock, right=0.7cm of fc1] (fc2) {};
\node[lbl, below=0.15cm of fc2] {FC 2\\(Dropout)};

%% ---- Output ----
\node[outblock, right=0.7cm of fc2] (out) {Softmax\\output};
\node[lbl, below=0.15cm of out] {Class\\probs.};

%% ---- Arrows ----
\draw[arrow] (img.east)    -- (conv10.west);
\draw[arrow] (conv12.east) -- (pool10.west);
\draw[arrow] (pool12.east) -- (conv20.west);
\draw[arrow] (conv23.east) -- (pool20.west);
\draw[arrow] (pool23.east) -- (conv30.west);
\draw[arrow] (conv34.east) -- node[above, font=\scriptsize]{Flatten} (fc1.west);
\draw[arrow] (fc1.east)    -- (fc2.west);
\draw[arrow] (fc2.east)    -- (out.west);

%% ---- Brace annotations above ----
\draw[decorate, decoration={brace, amplitude=5pt, raise=4pt}, thick, blue!60]
    (conv10.north west) -- (pool12.north east)
    node[midway, above=8pt, font=\scriptsize, blue!80] {Feature extraction stage 1};

\draw[decorate, decoration={brace, amplitude=5pt, raise=4pt}, thick, blue!60]
    (conv20.north west) -- (pool23.north east)
    node[midway, above=12pt, font=\scriptsize, blue!80] {Feature extraction stage 2};

\draw[decorate, decoration={brace, amplitude=5pt, raise=3pt}, thick, blue!60]
    (conv30.north west) -- (conv34.north east)
    node[midway, above=14pt, font=\scriptsize, blue!80] {Deep feature maps};

\draw[decorate, decoration={brace, amplitude=5pt, raise=3pt}, thick, orange!70]
    (fc1.north west) -- (fc2.north east)
    node[midway, above=10pt, font=\scriptsize, orange!90] {Classification head};

\end{tikzpicture}
 from chapter4.tex
% Requires: tikz, tikz libraries (shapes.geometric, arrows.meta, positioning, fit, backgrounds)

\begin{tikzpicture}[
    font=\small,
    >=Stealth,
    node distance=0.55cm and 0.9cm,
    % ---- block styles ----
    imgblock/.style={
        draw, fill=pink!40, rounded corners=2pt,
        minimum width=0.7cm, minimum height=1.6cm,
        align=center
    },
    convblock/.style={
        draw, fill=blue!25, rounded corners=2pt,
        minimum width=0.55cm, minimum height=1.4cm,
        align=center
    },
    poolblock/.style={
        draw, fill=green!30, rounded corners=2pt,
        minimum width=0.55cm, minimum height=1.0cm,
        align=center
    },
    fcblock/.style={
        draw, fill=orange!30, rounded corners=2pt,
        minimum width=0.55cm, minimum height=0.9cm,
        align=center
    },
    outblock/.style={
        draw, fill=red!25, rounded corners=4pt,
        minimum width=1.1cm, minimum height=0.7cm,
        align=center
    },
    lbl/.style={font=\scriptsize, align=center},
    arrow/.style={->, thick, gray!70}
]

%% ---- Input image ----
\node[imgblock] (img) {};
\node[lbl, below=0.15cm of img] {Input\\image\\patch};

%% ---- Conv block 1 (3 feature maps shown as stacked rectangles) ----
\foreach \i in {0,1,2}{
    \node[convblock, right=\i*0.12cm+0.9cm of img, yshift=\i*0.12cm] (conv1\i) {};
}
\node[lbl, below=0.15cm of conv10] {Conv 1\\(ReLU)};

%% ---- Pool 1 ----
\foreach \i in {0,1,2}{
    \node[poolblock, right=\i*0.12cm+0.55cm of conv10, yshift=\i*0.12cm] (pool1\i) {};
}
\node[lbl, below=0.15cm of pool10] {Pool 1\\(Max)};

%% ---- Conv block 2 ----
\foreach \i in {0,1,2,3}{
    \node[convblock, right=\i*0.12cm+0.55cm of pool10, yshift=\i*0.12cm] (conv2\i) {};
}
\node[lbl, below=0.15cm of conv20] {Conv 2\\(BN+ReLU)};

%% ---- Pool 2 ----
\foreach \i in {0,1,2,3}{
    \node[poolblock, right=\i*0.12cm+0.55cm of conv20, yshift=\i*0.12cm] (pool2\i) {};
}
\node[lbl, below=0.15cm of pool20] {Pool 2\\(Max)};

%% ---- Conv block 3 ----
\foreach \i in {0,...,4}{
    \node[convblock, right=\i*0.12cm+0.55cm of pool20, yshift=\i*0.12cm] (conv3\i) {};
}
\node[lbl, below=0.15cm of conv30] {Conv 3\\(BN+ReLU)};

%% ---- Flatten arrow label ----
\node[right=0.65cm of conv34, yshift=-0.3cm] (flat) {};

%% ---- FC layers ----
\node[fcblock, right=1.5cm of conv30] (fc1) {};
\node[lbl, below=0.15cm of fc1] {FC 1\\(Dropout)};

\node[fcblock, right=0.7cm of fc1] (fc2) {};
\node[lbl, below=0.15cm of fc2] {FC 2\\(Dropout)};

%% ---- Output ----
\node[outblock, right=0.7cm of fc2] (out) {Softmax\\output};
\node[lbl, below=0.15cm of out] {Class\\probs.};

%% ---- Arrows ----
\draw[arrow] (img.east)    -- (conv10.west);
\draw[arrow] (conv12.east) -- (pool10.west);
\draw[arrow] (pool12.east) -- (conv20.west);
\draw[arrow] (conv23.east) -- (pool20.west);
\draw[arrow] (pool23.east) -- (conv30.west);
\draw[arrow] (conv34.east) -- node[above, font=\scriptsize]{Flatten} (fc1.west);
\draw[arrow] (fc1.east)    -- (fc2.west);
\draw[arrow] (fc2.east)    -- (out.west);

%% ---- Brace annotations above ----
\draw[decorate, decoration={brace, amplitude=5pt, raise=4pt}, thick, blue!60]
    (conv10.north west) -- (pool12.north east)
    node[midway, above=8pt, font=\scriptsize, blue!80] {Feature extraction stage 1};

\draw[decorate, decoration={brace, amplitude=5pt, raise=4pt}, thick, blue!60]
    (conv20.north west) -- (pool23.north east)
    node[midway, above=12pt, font=\scriptsize, blue!80] {Feature extraction stage 2};

\draw[decorate, decoration={brace, amplitude=5pt, raise=3pt}, thick, blue!60]
    (conv30.north west) -- (conv34.north east)
    node[midway, above=14pt, font=\scriptsize, blue!80] {Deep feature maps};

\draw[decorate, decoration={brace, amplitude=5pt, raise=3pt}, thick, orange!70]
    (fc1.north west) -- (fc2.north east)
    node[midway, above=10pt, font=\scriptsize, orange!90] {Classification head};

\end{tikzpicture}
 from chapter4.tex
% Requires: tikz, tikz libraries (shapes.geometric, arrows.meta, positioning, fit, backgrounds)

\begin{tikzpicture}[
    font=\small,
    >=Stealth,
    node distance=0.55cm and 0.9cm,
    % ---- block styles ----
    imgblock/.style={
        draw, fill=pink!40, rounded corners=2pt,
        minimum width=0.7cm, minimum height=1.6cm,
        align=center
    },
    convblock/.style={
        draw, fill=blue!25, rounded corners=2pt,
        minimum width=0.55cm, minimum height=1.4cm,
        align=center
    },
    poolblock/.style={
        draw, fill=green!30, rounded corners=2pt,
        minimum width=0.55cm, minimum height=1.0cm,
        align=center
    },
    fcblock/.style={
        draw, fill=orange!30, rounded corners=2pt,
        minimum width=0.55cm, minimum height=0.9cm,
        align=center
    },
    outblock/.style={
        draw, fill=red!25, rounded corners=4pt,
        minimum width=1.1cm, minimum height=0.7cm,
        align=center
    },
    lbl/.style={font=\scriptsize, align=center},
    arrow/.style={->, thick, gray!70}
]

%% ---- Input image ----
\node[imgblock] (img) {};
\node[lbl, below=0.15cm of img] {Input\\image\\patch};

%% ---- Conv block 1 (3 feature maps shown as stacked rectangles) ----
\foreach \i in {0,1,2}{
    \node[convblock, right=\i*0.12cm+0.9cm of img, yshift=\i*0.12cm] (conv1\i) {};
}
\node[lbl, below=0.15cm of conv10] {Conv 1\\(ReLU)};

%% ---- Pool 1 ----
\foreach \i in {0,1,2}{
    \node[poolblock, right=\i*0.12cm+0.55cm of conv10, yshift=\i*0.12cm] (pool1\i) {};
}
\node[lbl, below=0.15cm of pool10] {Pool 1\\(Max)};

%% ---- Conv block 2 ----
\foreach \i in {0,1,2,3}{
    \node[convblock, right=\i*0.12cm+0.55cm of pool10, yshift=\i*0.12cm] (conv2\i) {};
}
\node[lbl, below=0.15cm of conv20] {Conv 2\\(BN+ReLU)};

%% ---- Pool 2 ----
\foreach \i in {0,1,2,3}{
    \node[poolblock, right=\i*0.12cm+0.55cm of conv20, yshift=\i*0.12cm] (pool2\i) {};
}
\node[lbl, below=0.15cm of pool20] {Pool 2\\(Max)};

%% ---- Conv block 3 ----
\foreach \i in {0,...,4}{
    \node[convblock, right=\i*0.12cm+0.55cm of pool20, yshift=\i*0.12cm] (conv3\i) {};
}
\node[lbl, below=0.15cm of conv30] {Conv 3\\(BN+ReLU)};

%% ---- Flatten arrow label ----
\node[right=0.65cm of conv34, yshift=-0.3cm] (flat) {};

%% ---- FC layers ----
\node[fcblock, right=1.5cm of conv30] (fc1) {};
\node[lbl, below=0.15cm of fc1] {FC 1\\(Dropout)};

\node[fcblock, right=0.7cm of fc1] (fc2) {};
\node[lbl, below=0.15cm of fc2] {FC 2\\(Dropout)};

%% ---- Output ----
\node[outblock, right=0.7cm of fc2] (out) {Softmax\\output};
\node[lbl, below=0.15cm of out] {Class\\probs.};

%% ---- Arrows ----
\draw[arrow] (img.east)    -- (conv10.west);
\draw[arrow] (conv12.east) -- (pool10.west);
\draw[arrow] (pool12.east) -- (conv20.west);
\draw[arrow] (conv23.east) -- (pool20.west);
\draw[arrow] (pool23.east) -- (conv30.west);
\draw[arrow] (conv34.east) -- node[above, font=\scriptsize]{Flatten} (fc1.west);
\draw[arrow] (fc1.east)    -- (fc2.west);
\draw[arrow] (fc2.east)    -- (out.west);

%% ---- Brace annotations above ----
\draw[decorate, decoration={brace, amplitude=5pt, raise=4pt}, thick, blue!60]
    (conv10.north west) -- (pool12.north east)
    node[midway, above=8pt, font=\scriptsize, blue!80] {Feature extraction stage 1};

\draw[decorate, decoration={brace, amplitude=5pt, raise=4pt}, thick, blue!60]
    (conv20.north west) -- (pool23.north east)
    node[midway, above=12pt, font=\scriptsize, blue!80] {Feature extraction stage 2};

\draw[decorate, decoration={brace, amplitude=5pt, raise=3pt}, thick, blue!60]
    (conv30.north west) -- (conv34.north east)
    node[midway, above=14pt, font=\scriptsize, blue!80] {Deep feature maps};

\draw[decorate, decoration={brace, amplitude=5pt, raise=3pt}, thick, orange!70]
    (fc1.north west) -- (fc2.north east)
    node[midway, above=10pt, font=\scriptsize, orange!90] {Classification head};

\end{tikzpicture}
